% Chapter 4 — link: names in the text become Wikidata items (kym_entities), and only those
% that matter to the meme are kept (kym_entity_curation).

\gmzoomin{link}
\newcommand{\hl}[1]{\colorbox{gm@col@link!20!gmSurface}{\strut #1}}
\note{Chapter 4: who and what is this page about? Linking names to Wikidata.}

\begin{frame}{Doge, step 3: who and what is the page about?}
  \vskip1mm
  {\setlength{\fboxsep}{1.2pt}%
  \begin{minipage}{\linewidth}\normalsize\color{gmInk}
    \hl{Doge} is a slang term for ``dog'' that is primarily associated with pictures of
    \hl{Shiba Inus} (nicknamed ``Shibe'') and internal monologue captions on \hl{Tumblr}. These
    photos may be photoshopped to change the dog's face or captioned with interior monologues in
    \hl{Comic Sans} font. The primary meme and iconography associated with Doge was the Shiba Inu
    named \hl{Kabosu}, whose photos taken by her owner, Atsuko Sato, in early 2010, went viral
    \dots
  \end{minipage}}
  \vskip7mm
  \begin{tikzpicture}[x=1cm, y=1cm,
      item/.style={rounded corners=1.2mm, draw=gm@col@link, line width=.8pt, fill=gmSurface,
                   inner sep=2.5pt, align=center, text width=2.25cm, minimum height=1.05cm}]
    \node[item] at (0,0)   {\footnotesize Doge\\[-.4mm]{\tiny\color{gmInk2}Internet meme}\\{\tiny\color{gmInk2}Q15613810}};
    \node[item] at (2.75,0){\footnotesize Shiba Inu\\[-.4mm]{\tiny\color{gmInk2}dog breed}\\{\tiny\color{gmInk2}Q39315}};
    \node[item] at (5.5,0) {\footnotesize Tumblr\\[-.4mm]{\tiny\color{gmInk2}microblogging platform}\\{\tiny\color{gmInk2}Q384060}};
    \node[item] at (8.25,0){\footnotesize Comic Sans MS\\[-.4mm]{\tiny\color{gmInk2}typeface}\\{\tiny\color{gmInk2}Q868492}};
    \node[item] at (11,0)  {\footnotesize Kabosu\\[-.4mm]{\tiny\color{gmInk2}Japanese dog, meme celebrity}\\{\tiny\color{gmInk2}Q23486479}};
  \end{tikzpicture}
  \vskip3mm
  {\small\color{gmInk2}Each name becomes a \textbf{Wikidata} item: a shared, worldwide identity with its own facts.}
  \note{
  \begin{itemize}
    \item Layman version: the computer underlines the names in the text and looks each one up in
      Wikidata --- the encyclopedia of facts behind Wikipedia --- where every thing has an ID.
    \item ``Shiba Inu'' becomes Q39315, the dog breed. Tumblr is Q384060. Kabosu, the real dog, has
      her own item. Comic Sans MS is the typeface.
    \item Why IDs matter: ``Shiba Inu'' in the Doge page and ``shiba'' in another page become the
      same thing, and the graph can connect them.
    \item Doge itself gets its item from its Know Your Meme ID, recorded in Wikidata --- that link
      is certain, score 1.0.
  \end{itemize}}
\end{frame}

\begin{frame}{Not every name matters}
  \begin{columns}[t]
    \column{.47\textwidth}
      {\bfseries\color{gm@col@link}\faCheck\enspace kept}\par\vskip2mm
      \small
      \begin{tabular}{@{}l@{\hspace{3mm}}l@{}}
        Shiba Inu      & {\color{gmInk2}a tag and the text agree}\\
        Kabosu         & {\color{gmInk2}a tag and the text agree}\\
        Tumblr, Reddit & {\color{gmInk2}a platform}\\
        Dogecoin       & {\color{gmInk2}a whole tag names it}\\
        monologue      & {\color{gmInk2}the judge: its format}\\
        Doge           & {\color{gmInk2}the meme's own item}\\
      \end{tabular}
    \column{.49\textwidth}
      {\bfseries\color{gmMuted}\faTimes\enspace dropped}\par\vskip2mm
      \small
      \begin{tabular}{@{}l@{\hspace{3mm}}l@{}}
        ``doges'' $\to$ the doges of Venice & {\color{gmInk2}wrong sense}\\
        ``Sato'' $\to$ slang for a stray dog & {\color{gmInk2}wrong sense}\\
        Donald Trump                       & {\color{gmInk2}incidental}\\
        Internet, internet meme            & {\color{gmInk2}incidental}\\
        animal, iconography, plan          & {\color{gmInk2}incidental}\\
      \end{tabular}
  \end{columns}
  \vskip6mm
  {\small\color{gmInk2}Doge: \NDogeMentions\ names found, \NDogeKept\ kept --- each with the reason it was kept.}
  \note{
  \begin{itemize}
    \item Recognising names is easy; deciding which ones \emph{matter to the meme} is the hard part.
    \item Real decisions on Doge. Kept: the breed, the dog, the platforms, the coin, the
      monologue format. Dropped: the tag ``doges'' linked to the doge, head of state of Venice and
      Genoa --- a wrong sense. ``Sato'' (Atsuko Sato, Kabosu's owner) linked to Puerto Rican slang
      for a mixed-breed dog --- a wrong sense, and a funny one for a dog meme. Donald Trump appears because of DOGE, the US department ---
      incidental to the meme.
    \item Every kept link carries its reason (rules first, the LLM judge for the rest), so anyone
      can see why it is in the graph.
  \end{itemize}}
\end{frame}

\begin{frame}{\NMentions\ names, \NItems\ things, \NKept\ kept}
  \begin{columns}[t]
    \column{.5\textwidth}
      {\scriptsize\color{gmInk2}kept, and why (\NKeptPct\%)}\par\vskip1mm
      % generated by scripts/figures.py — do not edit
\begin{tikzpicture}
\begin{axis}[gm hbar, scale only axis, width=.8\linewidth, height=4.3cm, xmax=66808, symbolic y coords={{a tag and the text agree},{the LLM judge},{the title agrees},{a platform},{a whole tag names it},{named by the title},{a meme format},{the meme's own item}},
  point meta=explicit symbolic, nodes near coords={\pgfplotspointmeta}]
\addplot[fill=gm@col@link] coordinates {(52194,{a tag and the text agree}) [52,194] (40783,{the LLM judge}) [40,783] (25685,{the title agrees}) [25,685] (25622,{a platform}) [25,622] (22428,{a whole tag names it}) [22,428] (17882,{named by the title}) [17,882] (13614,{a meme format}) [13,614] (4475,{the meme's own item}) [4,475]};
\end{axis}
\end{tikzpicture}

    \column{.46\textwidth}
      {\scriptsize\color{gmInk2}dropped, and by what}\par\vskip1mm
      % generated by scripts/figures.py — do not edit
\begin{tikzpicture}
\begin{axis}[gm hbar, scale only axis, width=.8\linewidth, height=2.0cm, xmax=101182, symbolic y coords={{the LLM judge},{too generic},{deny list: item},{deny list: class}},
  point meta=explicit symbolic, nodes near coords={\pgfplotspointmeta}]
\addplot[fill=gmMuted] coordinates {(67455,{the LLM judge}) [67,455] (21375,{too generic}) [21,375] (3454,{deny list: item}) [3,454] (3248,{deny list: class}) [3,248]};
\end{axis}
\end{tikzpicture}

      \vskip3mm
      \gmbig{\NLexEntities}{Wikidata items in our own lexicon, from the \NDumpGB\ GB dump}
  \end{columns}
  \note{
  \begin{itemize}
    \item 305,359 names found in 24,242 entries, pointing at 31,085 different Wikidata items.
    \item Curation kept 209,827 links (68.7\%) and dropped 95,532.
    \item Kept: mostly by rules --- a tag and the text agree, a platform, the title agrees --- and
      42,114 by the LLM judge. Dropped: mostly by the judge (67,455), then items too generic to say
      anything (21,375), then curated deny lists.
    \item The lexicon: 17.3 million Wikidata items with 22.6 million names, built from the full
      156 GB Wikidata dump (121.7 million lines), stored locally (3.3 GB).
  \end{itemize}}
\end{frame}

\begin{frame}{Linking like IMKG --- but at home, with evidence}
  \begin{columns}[t]
    \column{.47\textwidth}
      {\scriptsize\bfseries\color{gmMuted}IMKG}\par\vskip1mm
      \begin{itemize}\setlength{\itemsep}{2mm}\small
        \item DBpedia Spotlight, a web service, confidence $\ge$ 0.5
        \item then DBpedia $\to$ Wikidata
        \item a link, and nothing about it
      \end{itemize}
    \column{.49\textwidth}
      {\scriptsize\bfseries\color{gm@col@link}MemeAtlas}\par\vskip1mm
      \begin{itemize}\setlength{\itemsep}{2mm}\small
        \item language analysis \textbf{on our machine} (spaCy): names, noun phrases, whole tags
        \item a \textbf{lexicon of \NLexEntities\ items} built from the full Wikidata dump
        \item every link keeps its \textbf{score, method, and reason}
        \item the meme's own item from its \textbf{KYM ID} in Wikidata: score 1.0
      \end{itemize}
  \end{columns}
  \vskip6mm
  \begin{tikzpicture}
    \node[fill=gmTint, rounded corners=2mm, inner sep=3mm, text width=12.8cm, align=left] {
      \gmwhy{Why at home?}\enspace\small The same dump gives the same links, every time --- no quota, no
      service that changes its answers, no second knowledge base in between.};
  \end{tikzpicture}
  \note{
  \begin{itemize}
    \item IMKG used DBpedia Spotlight, a remote API, then mapped DBpedia to Wikidata. Remote,
      rate-limited, and its answers change over time.
    \item MemeAtlas: spaCy (en\_core\_web\_sm) finds candidate spans --- named entities, noun
      chunks, runs of proper nouns, the title, each whole tag --- and a local lexicon built from
      the dated Wikidata dump (wikidata-20260914) links them. Same dump, same answer.
    \item Every link records score, method, the NER label, and the reason curation kept it. In RDF
      these are RDF-star annotations on the link itself.
    \item The meme's own item: Wikidata has a property for Know Your Meme IDs (P13484, the slug);
      when it matches, the link is certain.
    \item IMKG's predicates are kept word for word: m4s:fromAbout, m4s:fromTags.
  \end{itemize}}
\end{frame}

\begin{frame}{Measuring changed the design}
  \begin{columns}[t]
    \column{.5\textwidth}
      \begin{itemize}\setlength{\itemsep}{2.6mm}
        \item \textbf{\FCurCommonNouns\%} of the About links came from common nouns:
          ``popularity'', ``man'', ``photograph''
        \item ``a series of videos'' $\to$ the \textbf{mathematical series}, \FCurSeries\ times
        \item So: a curated list of senses, then \textbf{curation} --- rules first, then an LLM judge
          asked twice, in two wordings
      \end{itemize}
    \column{.44\textwidth}
      {\scriptsize\color{gmInk2}the bar, set before measuring --- and the result}\par\vskip2mm
      \begin{tabular}{@{}lcc@{}}
        & \small bar & \small measured\\\midrule
        \small relevant among kept    & $\ge$\,\FCurBarKept & \textbf{\FCurKept}\\
        \small relevant among dropped & $\le$\,\FCurBarDropped & \textbf{\FCurDropped}\\
      \end{tabular}\par\vskip3mm
      {\scriptsize\color{gmMuted}projected over the corpus from 353 labelled links}
  \end{columns}
  \note{
  \begin{itemize}
    \item The first full linking run (282,914 links) was measured before anything else was built
      on it. 59\% of About links came from spans with no proper noun --- common-noun concepts.
    \item Worse, some hubs were wrong: KYM's stock phrase ``X is a series of \dots'' linked to the
      mathematical series (Q170198) 5,531 times; ``game'' to games in general, where KYM means
      video games.
    \item Fixes: a curated sense list (entity\_senses.yaml); then a curation stage: rules decide
      first (title, own item, platform, format, a tag and the text agree, a whole tag), an LLM
      (ministral-3:14b on the lab's GPUs) judges the rest with two differently worded prompts.
    \item The bar was set before measuring: of kept links $\ge$ 0.85 relevant, of dropped $\le$
      0.15. A first review failed on the dropped side (0.25), which produced curation 1.1.0; it
      now projects to 0.875 / 0.126. Caveat to say out loud: projected from 353 labelled items; the
      holdout review is the next step.
  \end{itemize}}
\end{frame}

\gmzoomout{link}
